\documentclass[10pt,a4paper]{amsart}
\usepackage[margin=19mm]{geometry}
\usepackage{amsmath,amssymb,mathtools}
\usepackage[scr=rsfs,cal=euler]{mathalfa}
\usepackage{xfrac}
\usepackage[unicode,hidelinks,bookmarks=false]{hyperref}
\newcommand{\ie}{i.e.\ }
\newcommand{\cf}{cf.\ }
\newcommand{\resp}{respectively\ }
\newcommand{\Subsection}{\subsection}
\providecommand{\nobreakdash}{\nobreak\hbox{-}}
\setlength{\emergencystretch}{2em}
\multlinegap0pt
\usepackage{mathtools}
\usepackage{amssymb}
\usepackage{xcolor}
\newtheorem{lem}{Lemma}
\DeclareMathOperator{\LL}{L}
\newcommand{\N}{\mathbb{N}}
\newcommand{\lin}{\operatorname{span}}
\newcommand{\mc}{\mathcal}
\newcommand{\md}{d}
\newcommand{\ov}{\overline}
\newcommand{\wt}{\widetilde}
\title{Crossed product functors associated to $\ell^p$-pseudofunctions}
\author{Jacek Krajczok, Ebrahim Samei, Timo Siebenand, Adam Skalski}
\date{Source: arXiv:2604.26345v1 (CC BY 4.0)}
\begin{document}
\maketitle

\noindent\emph{Source and licence.} Crossed product functors associated to $\ell^p$-pseudofunctions. Version arXiv:2604.26345v1 (CC BY 4.0), distributed by arXiv under the Creative Commons Attribution 4.0 International licence, http://creativecommons.org/licenses/by/4.0/, as recorded in the arXiv licence metadata for this exact version and shown on the abstract page https://arxiv.org/abs/2604.26345v1. Full licence notice: \texttt{licenses/arxiv-2604.26345-CC-BY-4.0.txt}.\par\smallskip

\noindent\textit{Editorial scope.} Counted unit: the article's lem lemma1 (source0.tex lines 684--692). The article's complete original proof of it is delivered with it (source0.tex lines 694--882, one \texttt{proof} environment of 189 lines). No mathematics is written, repaired or re-derived: the statement and the proof are the article's own lines, in the article's own notation, and no retained line is substituted or re-flowed.  Retained uncounted as the source-local context the proof needs: lines 676--680.  Nothing was omitted from any retained span, so the package carries no omission record.  No retained line contains a citation, so the wrapper carries no bibliography and no bibliography entry.  No retained line contains a figure environment, an image-inclusion command or drawing code, so no image is delivered and the package declares no asset dependency.  Editorial formatting only, in addition: an AMS \texttt{amsart} preamble that restates the article's own declarations and macros used by the retained lines, namely the environments \texttt{lem} on separate counters, together with the standard packages those lines need.  The retained environments print in this document as Lemma 1 in order of appearance, whereas the article prints the same items with the numbering of its own full document.  The complete native source of the article as distributed in the arXiv e-print for this exact version is bundled unchanged as \texttt{sources/199387/source0.tex}.
\begin{equation} \label{KKInequality}
\bigl(\int_{0}^{1}\bigl\|\sum_{k=1}^{n} x_k r_k(t)\bigr\|^q \md t\bigr)^{1/q}\le
K_{p,q}
\bigl(\int_{0}^{1}\bigl\|\sum_{k=1}^{n} x_k r_k(t)\bigr\|^p \md t\bigr)^{1/p}
\end{equation}

\begin{lem}\label{lemma1}
Let $1\le p<\infty$, $A$ be a $G$-C$^*$-algebra, $\pi$ a non-degenerate $*$-representation of $A$ on a Hilbert space $\mc{H}_\pi$ and $\sigma=\pi(\cdot)\otimes I_{\mc{K}}\colon A\rightarrow \mc{B}(\mc{H}\otimes \mc{K})$ its amplification for a non-zero Hilbert space $\mc{K}$. Then there exists a constant $C_p\ge 1$ depending only on $p$, such that
\begin{equation}\label{eq3}
\|(\widetilde{\pi}_p\rtimes \lambda^p)(f)\|_{\mc{B}(\ell^p(G;\mc{H}))}\le
\|(\widetilde{\sigma}_p\rtimes \lambda^p)(f)\|_{\mc{B}(\ell^p(G;\mc{H}\otimes \mc{K}))}\le
C_p \|(\widetilde{\pi}_p\rtimes \lambda^p)(f)\|_{\mc{B}(\ell^p(G;\mc{H}))}
\end{equation}
for $f\in\ell^1(G;A).$
\end{lem}

\begin{proof}
Since representations $\widetilde{\pi}_p\rtimes \lambda^p, \widetilde{\sigma}_p\rtimes \lambda^p$ are contractive, it is enough to prove \eqref{eq3} for a fixed finitely supported function $f\colon G\rightarrow A$.

The first inequality in \eqref{eq3} is straightforward: take $\xi\in \ell^p(G;\mc{H})$, choose a unit vector $\zeta\in\mc{K}$ and define $\xi_\zeta\colon G\ni g\mapsto \xi(g)\otimes \zeta\in \mc{H}\otimes\mc{K}$. Then $\|\xi_\zeta\|_{\ell^p(G;\mc{H}\otimes \mc{K})}=\|\xi\|_{\ell^p(G;\mc{H})}$ and
\begin{align*}
\|(&\wt{\pi}_p\rtimes\lambda^p)(f)\xi\|_{\ell^p(G;\mc{H})}=
\bigl(\sum_{h\in G}\bigl\|\sum_{g\in G}  \pi\bigl(\alpha_{h^{-1}}(f(g))\bigr)\xi(g^{-1} h)\bigr\|^p_{\mc{H}}\bigr)^{1/p}\\
&=
\bigl(\sum_{h\in G}\bigl\|\sum_{g\in G}  \pi\bigl(\alpha_{h^{-1}}(f(g))\bigr)\xi(g^{-1} h)\otimes \zeta\bigr\|^p_{\mc{H}\otimes \mc{K}}\bigr)^{1/p}=
\bigl(\sum_{h\in G}\bigl\|
((\wt{\sigma}_p\rtimes\lambda^p)(f) \xi_\zeta)(h)
\bigr\|^p_{\mc{H}\otimes \mc{K}}\bigr)^{1/p}\\
&=
\bigl\|
(\wt{\sigma}_p\rtimes\lambda^p)(f) \xi_\zeta\|_{\ell^p(G;\mc{H}\otimes \mc{K})} \le
\|(\wt{\sigma}_p\rtimes\lambda^p)(f)\| \|\xi\|_{\ell^p(G;\mc{H})},
\end{align*}
which proves the first inequality in \eqref{eq3}. Fix $\xi\in \ell^p(G;\mc{H}\otimes \mc{K})$ with norm $\le 1$. To establish the second inequality, we need to prove that
\begin{equation}\label{eq4}
\|(\widetilde{\sigma}_p\rtimes \lambda^p)(f)\xi\|_{\ell^p(G;\mc{H}\otimes \mc{K})}\le
C_p \|(\widetilde{\pi}_p\rtimes \lambda^p)(f) \|_{\ell^p(G;\mc{H})}
\end{equation}
with a constant $C_p$ depending only on $p$. Since $\xi$ is $p$-summable, there is a separable Hilbert space $\mc{K}_0\subseteq \mc{K}$ such that $\xi(h)\in \mc{H}\otimes \mc{K}_0$ for $h\in G$. Then we easily see that \eqref{eq4} is equivalent to
\begin{equation}\label{eq5}
\|(\widetilde{\sigma}_p\rtimes \lambda^p)(f)\xi\|_{\ell^p(G;\mc{H}\otimes \mc{K}_0)}\le
C_p \|(\widetilde{\pi}_p\rtimes \lambda^p)(f) \|_{\ell^p(G;\mc{H})},
\end{equation}
where by a slight abuse of the notation we write also $\sigma$ for the amplification of $\pi$ on $\mc{H}\otimes \mc{K}_0$. By passing via isometry, we can in fact assume that $$ \mc{K}_0=\mc{R}=\ov{\lin}\{r_n\mid n\in\N\}\subseteq \LL^2(\left[0,1\right]).$$ For $h\in G$, we write
\[
\xi(h)=
\sum_{n=1}^{\infty}  \eta_n(h)\otimes r_n,\quad
\textnormal{where}\;\;\;
\eta_n\colon G\ni h\mapsto(\mathrm{id}\otimes   r_n^*) (\xi(h))\in \mc{H}.
\]
The above series converges in the norm of $\mc{H}\otimes \mc{R}$ for every $h\in G$. Then also
\[
\bigl\|\xi-\sum_{n=1}^{N} \eta_n(\cdot)\otimes r_n\bigr\|_{\ell^p(G;\mc{H}\otimes\mc{R})}^p
=
\sum_{h\in G} \bigl\|
\sum_{n=N+1}^{\infty} \eta_n(h)\otimes r_n\bigr\|_{\mc{H}\otimes \mc{R}}^p \xrightarrow[N\to\infty]{}0
\]
by the bounded convergence theorem, where the uniform bound is given by
\[
\bigl\|
\sum_{n=N+1}^{\infty} \eta_n(h)\otimes r_n\bigr\|_{\mc{H}\otimes \mc{R}}^p=
\bigl(\sum_{n=N+1}^{\infty} \bigl\|\eta_n(h)\bigr\|_{\mc{H}}^{2 }\bigr)^{p/2}\le
\bigl(\sum_{n=1}^{\infty} \bigl\|\eta_n(h)\bigr\|_{\mc{H}}^{2}\bigr)^{p/2}=
\|\xi(h)\|^{p}.
\]
This proves that $\sum_{n=1}^{N} \eta_n(\cdot)\otimes r_n\xrightarrow[N\to\infty]{}\xi$ in $\ell^p(G;\mc{H}\otimes\mc{R})$. With these preliminaries we calculate the left hand side of \eqref{eq5}
\begin{align*}
\|(\wt{\sigma}_p\rtimes\lambda^p)(f) \xi\|_{\ell^p(G;\mc{H}\otimes \mc{R})}&=
\lim_{N\to\infty}
\bigl\|(\wt{\sigma}_p\rtimes\lambda^p)(f)
\bigl(
\sum_{n=1}^{N} \eta_n(\cdot)\otimes r_n
\bigr)
\bigr\|_{\ell^p(G;\mc{H}\otimes \mc{R})}
\\
&=
\lim_{N\to\infty}
\bigl(
\sum_{h\in G}
\bigl\|
\sum_{g\in G}
\bigl(
\wt{\sigma}_p(f(g))\lambda^p(g)
\bigl(\sum_{n=1}^{N}\eta_n(\cdot)\otimes r_n\bigr)
\bigr)(h)\bigr\|^p_{\mc{H}\otimes\mc{R}}\bigr)^{1/p}\\
&=
\lim_{N\to\infty}
\bigl(
\sum_{h\in G}
\bigl\|
\sum_{g\in G}
\sum_{n=1}^{N}
 \pi (\alpha_{h^{-1}}(f(g)))
\eta_n(g^{-1}h)\otimes r_n
\bigr\|^p_{\mc{H}\otimes\mc{R}}\bigr)^{1/p}.
\end{align*}
Note that the sum over $g\in G$ is finite. As $\mc{R}\subseteq \LL^2([0,1])$, we can calculate the last norm in $\mc{H}\otimes \LL^2([0,1])\simeq \LL^2([0,1]; \mc{H})$, using \eqref{KKInequality}:
\begin{align*}
\|(\wt{\sigma}_p\rtimes\lambda^p)(f) \xi&\|_{\ell^p(G;\mc{H}\otimes \mc{R})}\\
&= \lim_{N\to\infty}
\bigl(
\sum_{h\in G}
\bigl(
\int_{0}^{1}
\bigl\|
\sum_{g\in G}
\sum_{n=1}^{N}
 \pi (\alpha_{h^{-1}}(f(g)))
\eta_n(g^{-1}h)r_n(t)
\bigr\|^2_{\mc{H}}
\md t
\bigr)^{p/2}\bigr)^{1/p}\\
&\le
\limsup_{N\to\infty}
\bigl(
\sum_{h\in G}
\bigl(
K_{p,2}
\bigl[
\int_{0}^{1}
\bigl\|
\sum_{g\in G}
\sum_{n=1}^{N}
 \pi(\alpha_{h^{-1}}(f(g)))
\eta_n(g^{-1}h)r_n(t)
\bigr\|^p_{\mc{H}}
\md t
\bigr]^{1/p}\bigr)^{p}
\bigr)^{1/p}\\
&=
K_{p,2}
\limsup_{N\to\infty}
\bigl(
\int_{0}^{1}
\sum_{h\in G}
\bigl\|
\sum_{n=1}^{N} r_n(t)
((\wt{\pi}_p\rtimes\lambda^p)(f)\eta_n)(h) \bigr\|^p_{\mc{H}}
\md t
\bigr)^{1/p}.
\end{align*}
We have switched the order of integral and sum, since the integrand is non-negative and we can restrict to a countable subset of $G$. Then
\begin{align*}
\|(\wt{\sigma}_p\rtimes\lambda^p)(f) \xi&\|_{\ell^p(G;\mc{H}\otimes \mc{R})}\\
&\le
K_{p,2}
\limsup_{N\to\infty}
\bigl(
\int_{0}^{1}
\bigl\|
\sum_{n=1}^{N} r_n(t)
(\wt{\pi}_p\rtimes\lambda^p)(f)\eta_n \bigr\|^p_{\ell^p(G;\mc{H})}
\md t
\bigr)^{1/p}\\
&\le
K_{p,2}
\|(\wt{\pi}_p\rtimes\lambda^p)(f)\|
\limsup_{N\to\infty}
\bigl(
\int_{0}^{1}
\bigl\|
\sum_{n=1}^{N} r_n(t)
 \eta_n \bigr\|^p_{\ell^p(G;\mc{H})}
\md t
\bigr)^{1/p}\\
&=
K_{p,2}
\|(\wt{\pi}_p\rtimes\lambda^p)(f)\|
\limsup_{N\to\infty}
\bigl(
\sum_{h\in G}
\int_{0}^{1}
\bigl\|
\sum_{n=1}^{N} r_n(t)
 \eta_n(h) \bigr\|^p_{ \mc{H}}
\md t
\bigr)^{1/p}\\
&\le
K_{p,2}
\|(\wt{\pi}_p\rtimes\lambda^p)(f)\|
\limsup_{N\to\infty}
\bigl(
\sum_{h\in G}
\bigl(
K_{2,p}\bigl[\int_{0}^{1}
\bigl\|
\sum_{n=1}^{N} r_n(t)
 \eta_n(h) \bigr\|^2_{ \mc{H}}
\md t\bigr]^{1/2}\bigr)^{p}
\bigr)^{1/p}\\
&=
K_{p,2} K_{2,p}
\|(\wt{\pi}_p\rtimes\lambda^p)(f)\|
\limsup_{N\to\infty}
\bigl\|
\sum_{n=1}^{N}
 \eta_n(\cdot)\otimes r_n \bigr\|_{\ell^p(G; \mc{H}\otimes \mc{R})}\\
 &=
K_{p,2} K_{2,p}
\|(\wt{\pi}_p\rtimes\lambda^p)(f)\|
\|
\xi \|_{\ell^p(G; \mc{H}\otimes \mc{R})}.
\end{align*}
This ends the proof with $C_p=K_{p,2} K_{2,p}$.
\end{proof}

\end{document}
